\chapter{Introduction}
\label{ch:intro}

% The subsections below match the requirements you provided (1.1–1.8).
\section{Introduction}
\noindent This chapter introduces the research topic and frames the significance of designing an accessible AI-based mental health monitoring system for low-resource settings. The scaffold below lists required subsections; expand each item with original writing, data, and citations.

\section{Background and Motivation}
\begin{itemize}
  \item Context: mental health access gaps in low-resource settings (brief summary).
  \item Motivation: value of USSD for offline/low-bandwidth access; AI for early risk detection.
\end{itemize}

\section{Research Problem Statement}
\noindent Concisely state the problem your system addresses, who is affected, and why current solutions are insufficient.

\section{Research Questions}
\begin{itemize}
  \item Example: How can an accessible USSD + web system detect early mental health risk signals in low-resource settings?
  \item Add 3–6 specific research questions tied to evaluation metrics.
\end{itemize}

\section{Aims and Objectives}
\noindent Provide the main aim and 4–6 specific objectives (design, implementation, evaluation, deployment considerations).

\section{Delineation and Assumptions}
\noindent Scope (what is in/out), assumptions about connectivity, users, datasets, and ethical considerations.

\section{Research Methodology (Brief)}
\noindent One-paragraph summary of chosen methods (mixed methods: system development + experiments + user testing).

\section{Organisation of the Thesis}
\noindent Short paragraph summarising chapter contents and how they link.

% Guidance note
\vspace{1em}
\textbf{NOTE:} This file contains a scaffold and brief example text. I can expand any subsection into a full draft when you confirm the citation style and provide deadlines.
